\documentclass[11pt]{article}
\usepackage[a4paper,margin=25mm]{geometry}
\usepackage[T1]{fontenc}
\usepackage{lmodern,amsmath,amssymb,booktabs,hyperref}
\title{A finite-bandwidth diagnostic for the normal-population experiment}
\author{Supplementary Material S1}
\date{R computation, 25 October 2026}
\begin{document}
\maketitle

This Gaussian calculation examines why the leading weak-limit variance can
be inaccurate at the bandwidths used in the normal-population experiment.
It combines Gaussian integration with a first-order expansion of the
estimated cutpoints and compares the resulting linear approximation with
the full R Monte Carlo results. It does not give the exact finite-sample
variance of the full nonlinear estimator, select a bandwidth, or correct
interval coverage.

\section*{Linear diagnostic}
At the standard normal population, let $s=\Phi^{-1}(3/4)$,
$d=\varphi(s)$, $x_1=-s$, $x_2=s$, and $g=-1/[2(2d)^3]$. The MAD
influence function is
\[
 \phi_S(x)=\frac{1/2-1\{|x|\le s\}}{2d}.
\]
Write $K_h(y)=h^{-1}\varphi(y/h)$ and $\varphi_\sigma$ for the density
of $\mathrm N(0,\sigma^2)$. Gaussian convolution gives
\[
 \mu_h(x)=E K_h(x-X)=\varphi_{\sqrt{1+h^2}}(x),
\]
and the exact covariance of the two fixed-location contributions is
\[
 C_{ij}(h)=\frac{\exp\{-(x_i-x_j)^2/(4h^2)\}}{2\sqrt\pi h}
 \varphi_{\sqrt{1+h^2/2}}\!\left(\frac{x_i+x_j}{2}\right)
 -\mu_h(x_i)\mu_h(x_j).
\]
The product of the two Gaussian kernels is proportional to a Gaussian
kernel centred at $(x_i+x_j)/2$, with variance $h^2/2$. Integrating that
product against the standard normal density gives the displayed expression.

Using the first-order cutpoint displacement with the population derivatives
$\varphi'(-s)=sd$ and $\varphi'(s)=-sd$ gives the centred linear
contribution
\[
 W_h(X)=g\{K_h(-s-X)+K_h(s-X)-2\mu_h(s)\}
        -2sdg\,\phi_S(X).
\]
The median displacement cancels in the sum of the two endpoint
contributions. The MAD displacement remains. This expression combines
exact finite-bandwidth kernel covariance with a first-order cutpoint
expansion; it is not an exact finite-bandwidth influence representation
of the full estimator.

The limiting variance constant is
\[
 \mathcal A_F=\frac{2g^2d}{2\sqrt\pi}.
\]
At the deterministic bandwidth used in the study,
$h=\kappa_{\mathrm N}s n^{-1/5}=n^{-1/5}$, where $\kappa_{\mathrm N}=1/s$, numerical integration gives
the following standard deviations for the scaled linear contribution.
\begin{center}
\begin{tabular}{rrr}
\toprule
Sample size & $\{h\operatorname{Var}(W_h)/\mathcal A_F\}^{1/2}$
            & Full-estimator standardized SD \\
\midrule
200 & 0.800 & 0.893 \\
1000 & 0.826 & 0.873 \\
5000 & 0.866 & 0.888 \\
\bottomrule
\end{tabular}
\end{center}
The integrations are split at $-s$ and $s$ to handle the jumps in the
influence function. The largest absolute numerical value of $E W_h$ in
these three R calculations is below $3\times10^{-16}$. The rightmost
column gives the standardized full-estimator SDs from manuscript Table~6,
using 3000 simulated datasets per row. The comparison at
$n=5000$ supports the explanation that lower-order covariance and
cutpoint terms remain relevant. Nonlinear transformation and other
higher-order terms are not accounted for.

\section*{Reproduction}
From the replication-workspace root, run
\begin{verbatim}
Rscript computation/R/run.R diagnostic --mode full
Rscript computation/R/run.R clt --mode full
\end{verbatim}
The first command writes deterministic Gaussian integrations to
\path{computation/generated_r/full/diagnostic/}. The second generates
the six normal and lognormal variance-limit designs in
\path{computation/generated_r/full/clt/}. Their configuration records
retain the RNG definition, master seed, stream keys and source hashes.
The diagnostic uses native R integration and no random draws. Both calculations were executed for this release.
\end{document}
